% Gauss, Werke X/1 (1917), native90–92, printed83–85.
\SourceLeaf{090}\ArithmeticHead{83}{R}{ZERLEGUNGEN EINER PRIMZAHL IN DREI QUADRATE.}
\begin{center}\textit{Additamentum:}\par\medskip
Numerorum non primorum in terna quadrata distributio.\end{center}
\begin{center}\begin{tabular}{@{}rl@{}}
\textbf{10}&9, 1 $\vert$\\
\textbf{18}&9, 9 — 16, 1, 1 $\vert$\\
\textbf{21}&16, 4, 1 $\vert$\\
\textbf{22}&9, 9, 4 $\vert$\\
\textbf{25}&16, 9 $\vert$\\
\textbf{26}&25, 1 — 16, 9, 1 $\vert$\\
\textbf{27}&9, 9, 9 — 25, 1, 1 $\vert$\\
\textbf{30}&25, 4, 1 $\vert$\\
\textbf{33}&25, 4, 4 — 16, 16, 1 $\vert$\\
\textbf{34}&25, 9 — 16, 9, 9 $\vert$\\
\textbf{35}&25, 9, 1 $\vert$\\
\textbf{38}&36, 1, 1 — 25, 9, 4 $\vert$\\
\textbf{42}&25, 16, 1 $\vert$\\
\textbf{45}&36, 9 — 25, 16, 4 $\vert$\\
\textbf{46}&36, 9, 1 $\vert$\\
\textbf{49}&36, 9, 4 $\vert$\\
\textbf{50}&49, 1 — 25, 25 — 25, 16, 9 $\vert$\\
\textbf{51}&49, 1, 1 — 25, 25, 1 $\vert$\\
\textbf{54}&49, 4, 1 — 36, 9, 9 $\vert$\\
\textbf{57}&49, 4, 4 — 25, 16, 16 $\vert$\\
\textbf{58}&49, 9 $\vert$\\
\textbf{62}&49, 9, 4 — 36, 25, 1 $\vert$\\
\textbf{65}&64, 1 — 49, 16 — 36, 25, 4 $\vert$\\
\textbf{66}&64, 1, 1 — 49, 16, 1 — 25, 25, 16 $\vert$
\end{tabular}\end{center}
\begin{center}\textit{Selecta}\end{center}
\begin{center}\begin{tabular}{@{}rl@{}}
\textbf{205}&196, 9 — 169, 36 — 144, 36, 25 $\vert$\\
\textbf{221}&196, 25 — 196, 16, 9 — 169, 36, 16 — 121, 100 — 121, 64, 36 $\vert$
\end{tabular}\end{center}
\par\bigskip\begin{center}\textbf{BEMERKUNGEN.}\end{center}
\begingroup\small
\makeatletter\tagsleft@true\let\veqno\@@leqno\makeatother
\par\indent Diese Aufzeichnung von \textsc{Gauss} ist wohl als ein früher Anlauf zu dem Ziele zu betrachten, welches
er mit seinem berühmten Satze im art. 291 der \textit{Disquisitiones Arithmeticae} erreicht hat. Nach der Ausdrucksweise jener Zeit heißt Teiler von $x^2+Ay^2$ jede Primzahl $\Delta$, von der ein Vielfaches durch die Form
darstellbar, für welche also $M\Delta=x^2+Ay^2$ ist, allgemeiner eine quadratische Form $\Delta x^2+2qxy+ry^2$
mit der Determinante $-A$, deren erster Koeffizient gleich $\Delta$ und durch welche dann auch der Multiplikator $M$ darstellbar ist. Ist $\Delta=4n+1$, so folgt wegen $\left(\dfrac{-A}{\Delta}\right)=1$ auch $\left(\dfrac\Delta A\right)=1$ und
\[
(-1)^{\frac{\Delta-1}{2}}=1,
\]
\noindent die Form gehört also zum Hauptgeschlecht. Bei solcher Deutung würde der \textsc{Gauss}sche Satz also lauten:
\textls[100]{Wenn }$A$\textls[100]{ Primzahl ist, so gibt es soviel Formen} (Klassen) \textls[100]{des Hauptgeschlechts, als es verschiedene Arten gibt, }$A$\textls[100]{ in drei Quadrate zu zerlegen}, eine Aussage, die nicht ganz zutreffend wäre.

Bei dem Versuche, den Beweisgang aufzuklären, werde auch $A=4n+1$ vorausgesetzt.

\textls[100]{Zu 1.} Ist $A=a^2+b^2+c^2$ und $a^2+b^2\equiv1\ (\mathrm{mod.}\ 4)$, so wird
\[
x^2+Ay^2=x^2+(a^2+b^2+c^2)y^2
\]
\noindent für $x=ac$, $y=b$ resp. $x=bc$, $y=a$ gleich
\[
(a^2+b^2)(b^2+c^2)\ \text{resp.}\ (a^2+b^2)(a^2+c^2),
\]
\par\vfill\noindent\hfill11*\par\clearpage
\SourceLeaf{091}\ArithmeticHead{84}{L}{ARITHMETIK.\quad NACHLASS.}
\noindent zu jeder Zerlegung von $A$ gehört also ein Teiler $a^2+b^2\equiv1\ (\mathrm{mod.}\ 4)$ und eine Form
\[
(a^2+b^2)x^2+2qxy+ry^2
\]
\noindent mit der Determinante $-A$, durch welche auch die andern Teiler $b^2+c^2$, $a^2+c^2$, darstellbar sind; diese
drei \textls[100]{verbundenen} oder \textls[100]{zusammenhängenden} Teiler oder Multiplikatoren geben also nur eine Form.
\par\noindent\hspace*{\parindent}\textls[100]{Zu 2.} Ist $\Delta$ Primteiler von der Form $4n+1$ von $x^2+Ay^2$ d. h.
\[
M\mathpunct{.}\Delta=x^2+Ay^2,
\]
\noindent so ist wegen $\left(\dfrac{-A}{\Delta}\right)=1$ auch $\left(\dfrac{-\Delta}{A}\right)=1$ also auch
\[
N\mathpunct{.}A=x^2+\Delta y^2.
\]
\noindent Gilt nun die Behauptung in 2. für die Primzahl $\Delta$, so entspricht dem Theiler $A=4n+1$ eine Zerlegung
\begin{equation*}\Delta=m^2+n^2+p^2\tag*{(1)}\end{equation*}
\noindent und der Multiplikator $N=m^2+n^2$. Demnach findet sich
\begin{equation*}(m^2+n^2)A=x^2+(m^2+n^2+p^2)y^2\tag*{(2)}\end{equation*}
\noindent oder
\[
A=\frac{x^2+p^2y^2}{m^2+n^2}+y^2,
\]
\noindent also eine Zerlegung von $A$ in drei Quadrate
\begin{equation*}A=a^2+b^2+c^2,\tag*{(3)}\end{equation*}
\noindent wenn man schreibt $y=c$ und
\[
x^2+p^2c^2=(m^2+n^2)(a^2+b^2)=(an-bm)^2+(am+bn)^2
\]
\noindent d. i.
\[
x^2=(an-bm)^2,\quad p^2c^2=(am+bn)^2,
\]
\noindent woraus
\[
(am+bn+cp)(am+bn-cp)=0
\]
\noindent also bei passender Bezeichnung der Basen der Quadrate und Wahl ihrer Vorzeichen
\begin{equation*}am+bn+cp=0\tag*{(4)}\end{equation*}
\noindent hervorgeht. Ferner nimmt die Gleichung (2) die Gestalt an
\begin{equation*}(m^2+n^2)(a^2+b^2+c^2)=(an-bm)^2+(m^2+n^2+p^2)c^2.\quad\text{—}\tag*{(5)}\end{equation*}
\noindent Wie jedoch aus dieser Betrachtung eine zum Beweise der Behauptung in 2. führende »Depressionsmethode«
zu entnehmen ist, bleibt unklar.
\par\noindent\hspace*{\parindent}\textls[100]{Zu 4.} Wären $A=a^2+b^2+c^2$ mit $a^2+b^2\equiv1\ (\mathrm{mod.}\ 4)$ und $A=a'^2+b'^2+c'^2$ zwei Zerlegungen
von $A$, denen ein- und derselbe Primteiler $\Delta=4n+1$ zukäme, so bestände einerseits eine Gleichung
\[
(a^2+b^2)\Delta=x^2+Ay^2=x^2+(a'^2+b'^2+c'^2)y^2,
\]
\noindent andererseits müßte deshalb $\Delta$ nach dem Satze in 2. die Summe zweier der drei Quadrate $a'^2,b'^2,c'^2$, etwa
\[
\Delta=a'^2+c'^2
\]
\noindent sein. Aus
\[
(a^2+b^2)\Delta=x^2+Ay^2=x^2+(a^2+b^2+c^2)y^2
\]
\clearpage
\SourceLeaf{092}\ArithmeticHead{85}{R}{ZERLEGUNGEN EINER PRIMZAHL IN DREI QUADRATE.}
\noindent folgt aber (siehe 2.)
\[
\Delta=m^2+n^2+p^2,\quad am+bn+cp=0
\]
\noindent und die Gleichung (5) oder
\[
(m^2+n^2)A=x^2+\Delta y^2.
\]
\noindent Da
\[
A=b'^2+(a'^2+c'^2)=b'^2+\Delta\,1^2
\]
\noindent und $A$ Primzahl ist, so erschließt man hieraus
\[
m^2+n^2=r^2+\Delta s^2=r^2+(m^2+n^2+p^2)s^2
\]
\noindent was $s=0$ erfordert. Die Zerlegung von $\Delta$ würde also
\[
\Delta=a'^2+c'^2=r^2+p^2,
\]
\noindent dem entsprechend die Gleichung (4) die einfachere Gestalt
\[
ar+cp=0\ \text{d. i.}\ aa'+cc'=0
\]
\noindent annähme. Dies gibt die Beziehungen $a=\mu t$, $c=\mu u$, $a'=\nu t$, $c'=\nu u$, also wenn $t^2+u^2=S$ gesetzt wird, für $A$ die zwei Zerlegungen
\[
A=b^2+S\mu^2,\quad A=b'^2+S\nu^2,
\]
\noindent die mit einander identisch sein müssen, da $A$ Primzahl ist. So erhielte man $b=b'$, $\mu=\nu$ und die beiden
als verschieden vorausgesetzten Zerlegungen von $A$ würden einander gleich.
\begin{flushright}\textsc{Bachmann.}\end{flushright}
\ClosingRule\clearpage\endgroup
